% Zdroj: PDF priklady-podzim-2023.pdf, fyzická str. 18. Značka: specializace UI-SU, UI-ZPJ. Zařazeno: S1 (CSP / backtracking / arc consistency).
\section{Problém splňování podmínek (CSP)}
\szzotazka{podzim 2023}{28}{Problém splňování podmínek (CSP, backtracking, arc consistency)}

\noindent\textbf{Zadání.}\par
Definujte problém splňování podmínek (CSP z anglického Constraint Satisfaction Problem). Popište prohledávání s navracením (backtracking) jako způsob řešení CPS. Vysvětlete pojmy ,,hranová konzistence`` (arc consistency), ,,kontrola dopředu`` (forward checking) a ,,pohled dopředu`` (look ahead).

Pomocí CSP modelujte problém Sudoku, tj. problém umístění čísel $1, \dots, 9$ do mřížky $9 \times 9$ takovým způsobem, že se v žádném sloupci, řádku ani boxu velikosti $3 \times 3$ žádná cifra neopakuje více než jednou. Předpokládejte, že některá čísla jsou už v mřížce zadaná.

\medskip
\noindent\textbf{Náčrt řešení.}\par
\begin{itemize}[nosep]
  \item Definice CSP -- trojice (proměnné, domény, podmínky), obecně může být pro každou proměnnou jiná doména, podmínek je konečně mnoho. Podmínka je relace mezi proměnnými, podmnožina Kartézského součinu domén.
  \item Popis backtracking, arc consistency, forward checking a look ahead.
  \begin{itemize}[nosep]
    \item \textit{backtracking} -- přiřadíme hodnotu nějaké proměnné a opakujeme dokud nedostaneme validní řešení nebo spor. V případě sporu backtrackujeme a učiníme jiné rozhodnutí.
    \item \textit{arc consistency} -- podmínka (arc, pokud máme pouze binární podmínky) je konzistentní, pokud pro každou hodnotu z domény existuje hodnota pro ostatní proměnné, které tuto podmínku splní. Hodnoty pro které neexistuje ohodnocení ostatních proměnných můžeme z domény vyloučit.
    \item \textit{forward checking} -- po ohodnocení proměnné $x$ zkontrolujeme ostatní proměnné, které se vyskytují v podmínkách spolu s $x$. Odstraníme nekonzistentní hodnoty.
    \item \textit{look ahead} -- podobné jako forward check, ale kontrolujeme postupně všechny podmínky, dokud se nějaké proměnné mění doména. Tj. jde o silnější techniku než forward check. Obdobně můžeme říct, že po ohodnocení proměnné se problém udělá arc consistent.
  \end{itemize}
  \item Model sudoku -- každé políčko je jedna proměnná, každá proměnná má doménu $1, \dots, 9$. Podmínky jsou ,,All Different`` (případně nerovnost po dvojicích) na vhodné podmnožiny proměnných (tj, sloupce, řádky, 3x3 políčka). Typicky máme v sudoku už zadaná nějaká čísla, ta můžeme přiřadit pomocí domén nebo pomocí podmínek.
\end{itemize}
